% Texto para el Related Work (Chapter I, sección 1.2), antes del estado del arte.
% Pegar junto con docs/prisma_arxiv.tex (el diagrama de la figura).

\subsection{Literature search}

\textbf{Review question.} Has the effect of data treatments and of structural modifications of the VAE on the generation quality of latent diffusion models been studied for small, homogeneous, and multimodal datasets?

\textbf{Source and date.} The arXiv API was queried on 11 October 2026, restricted to submissions from 2022 onward. All strings combine a latent-diffusion term (L) with a term for the autoencoder (A), plus the term specific to each string:
\begin{itemize}
  \item L = \texttt{all:"latent diffusion" OR abs:"latent diffusion model"}
  \item A = \texttt{abs:autoencoder OR abs:VAE OR abs:tokenizer}
\end{itemize}

\begin{table}[h]
\centering
\small
\begin{tabular}{p{0.06\textwidth}p{0.62\textwidth}p{0.18\textwidth}}
\hline
\textbf{ID} & \textbf{Search string (in addition to L)} & \textbf{Records} \\
\hline
S1 & A \texttt{AND (abs:augmentation OR abs:"data augmentation")} & 30 \\
S2 & A \texttt{AND (abs:adversarial OR abs:"KL" OR abs:regularization OR abs:"latent channels")} & 73 \\
S3 & A \texttt{AND (abs:aerial OR abs:agricultural OR abs:crop OR abs:plant OR abs:UAV OR abs:"remote sensing")} & 9 \\
S4 & \texttt{(abs:masked OR abs:masking) AND} A & 47 \\
S5 & \texttt{abs:depth AND (abs:RGB OR abs:appearance) AND (abs:joint OR abs:jointly OR abs:"shared latent")} & 7 \\
\hline
\end{tabular}
\caption{Search strings and number of records returned by arXiv. A first version of S1 and S2 that also required a small-data term was too restrictive (it returned an error and one record, respectively) and was replaced by the strings above.}
\end{table}

\textbf{Screening criteria.} Records were screened by title and abstract with three criteria:
\begin{enumerate}
  \item \textbf{C1:} the work varies something in the autoencoder or in its training data (augmentation, masking, adversarial loss, KL or regularization, number of latent channels, or fine-tuning).
  \item \textbf{C2:} the work measures generation (FID, KID, generation or sample quality, image generation or synthesis).
  \item \textbf{C3:} the work addresses few or homogeneous data (small or limited datasets, few images, scarcity, few-shot, data efficiency) or an aerial, agricultural, or medical domain.
\end{enumerate}
The screening was automatic: each criterion was a set of keywords applied to the title and abstract, and the abstracts of the records that met C1, C2, and C3 were then read by the author. A work is considered to study the review question when it meets the three criteria and examines how the data or the structure of the autoencoder change generation.

\textbf{Result.} Figure~\ref{fig:prisma} summarizes the flow. Of 151 unique records, 34 met C1 and C2, and 4 of them also met C3 or the domain; the four are medical applications that use a VAE to synthesize data and do not study the effect of its design or of the data treatments. No record studied the review question.

\textbf{Limitations.} Only arXiv was searched, the screening relied on keywords, and the search was not repeated by a second reviewer; it is a first systematic search and not a full systematic review.
